\begin{table}[htbp]
    \centering

    % Número de tabla
    \textbf{Tabla \refstepcounter{table}\thetable}\label{tab:ejemplo-estructura}

    \vspace{0.5\baselineskip}

    % Título de la tabla
    \textit{El título debe ser breve y descriptivo}

    \vspace{0.5\baselineskip}

    \begin{tabular}{ll}
        \toprule

        Columna uno        & Columna dos        \\

        \midrule

        Número de la tabla & Título de la tabla \\

        Encabezado         & Datos de la tabla  \\

        Datos de la tabla  & Datos de la tabla  \\

        Datos de la tabla  & Datos de la tabla  \\

        Datos de la tabla  & Datos de la tabla  \\

        Datos de la tabla  & Datos de la tabla  \\

        \bottomrule
    \end{tabular}

    \vspace{0.5\baselineskip}

    \begin{flushleft}
        \textit{Nota.}
        Se hace una descripción del contenido de la tabla en
        función de la información que se esté exponiendo.
    \end{flushleft}

\end{table}